% !Mode:: "TeX:UTF-8"

\chapter{数据挖掘方法}

\section{描述性统计与准确率分析}

首先按模型、Prompt 层级、干扰条件和场景临界性统计准确率。该部分用于回答“哪个模型在哪类描述上更稳定”这一基础问题，并为后续可视化提供数据。

\section{层级一致性分析}

层级一致性用于衡量同一场景从 L1 到 L3 描述变化时，模型判断是否保持一致。如果模型在 L3 形式化描述中反而更容易出错，则说明形式化符号并不必然带来更强空间建模能力。

\section{噪声翻转率分析}

噪声翻转率衡量加入无关信息后，模型相对无噪声版本的最终判断是否发生变化。该指标可以识别模型是否过度关注颜色、噪声、光照和情绪压力等非几何因素。

\section{决策树与聚类分析}

决策树用于解释哪些特征最能区分正确与错误回复；K-Means 聚类用于根据准确率、一致性、公式使用率、坐标系使用率、噪声翻转率等指标划分模型行为类型。例如，某些模型可能表现为“符号逻辑型但计算链不稳定”，另一些模型可能表现为“自然语言直觉型但参数敏感性较低”。

\section{FP-Growth 关联规则挖掘}

FP-Growth 用于挖掘“条件组合 $\rightarrow$ 结果模式”的关联关系 \cite{han2000mining}。例如，可以分析如下规则：

\begin{equation}
\{level=L3,\ has\_noise=1\} \Rightarrow error\_category=calculation\_collapse
\label{eq:rule}
\end{equation}

这类规则有助于从大量回复中提炼可解释的失效路径，是本项目体现数据挖掘价值的重要部分。
